\begin{proof}[Proof of \cref{obj:lem:jackson-boundary-adaptive}]
\begin{enumerate}
\item Set
\[
\mathcal R=\prod_{\zeta\in\mathcal Z}
[c^0_\zeta-R_\zeta,c^0_\zeta+R_\zeta],
\qquad
\Delta(\mathbf w,\mathbf w')
=\sum_{\zeta\in\mathcal Z}|w_\zeta-w'_\zeta|,
\qquad
\Lambda_\epsilon=3(1+\epsilon^{-1})+1
\]
for \(\mathbf w,\mathbf w'\in\mathcal R\). The rectangle assumption makes both vectors nonnegative. We first establish
\[
|f_\lambda(\mathbf w)-f_\lambda(\mathbf w')|
\leq\Lambda_\epsilon\Delta(\mathbf w,\mathbf w').
\]

For each arm \(a\in\{0,1\}\), the contribution in \cref{obj:def:dual-process} satisfies
\[
|g_a(\mathbf w)-g_a(\mathbf w')|
\leq(1+\epsilon^{-1})\Delta(\mathbf w,\mathbf w').
\]
Here is the scalar estimate behind this bound. At zero total mass, every coordinate vanishes. We evaluate the quotient defining \(g_a\) in \cref{obj:def:dual-process} with the convention \(0/0=0\), so \(g_a(\mathbf 0)=0\). Since \(0\leq g_a(\mathbf w)\leq s(\mathbf w)\), and likewise for \(\mathbf w'\), the claim follows whenever either vector has zero total mass. For positive total masses, the defining maximum gives \(g_a(\mathbf w)=w_{a1}/\epsilon\) when \(s_a(\mathbf w)\leq\epsilon s(\mathbf w)\), and \(g_a(\mathbf w)=s(\mathbf w)w_{a1}/s_a(\mathbf w)\) in the other regime. When both vectors are in the first regime, their difference is bounded by \(\epsilon^{-1}\Delta(\mathbf w,\mathbf w')\). When both are in the second, \(0\leq w'_{a1}/s_a(\mathbf w')\leq1\) and
\[
\left|
\frac{w_{a1}}{s_a(\mathbf w)}
-\frac{w'_{a1}}{s_a(\mathbf w')}
\right|
\leq
\frac{|w_{a1}-w'_{a1}|+|w_{a0}-w'_{a0}|}
{s_a(\mathbf w)}.
\]
Expanding the difference of the two products, using
\(|s(\mathbf w)-s(\mathbf w')|\leq\Delta(\mathbf w,\mathbf w')\) and
\(s(\mathbf w)/s_a(\mathbf w)\leq\epsilon^{-1}\), proves the displayed arm bound.

For the mixed regime in which \(\mathbf w\) is in the first regime and \(\mathbf w'\) in the second, the difference equals
\[
\frac{w_{a1}-w'_{a1}}{\epsilon}
+
\frac{w'_{a1}}{s_a(\mathbf w')}
\frac{s_a(\mathbf w')-\epsilon s(\mathbf w')}{\epsilon}.
\]
The second term is nonnegative. Its ratio lies in \([0,1]\), and the two regime inequalities imply
\[
0\leq s_a(\mathbf w')-\epsilon s(\mathbf w')
\leq s_a(\mathbf w')-s_a(\mathbf w)
       -\epsilon\{s(\mathbf w')-s(\mathbf w)\}.
\]
Consequently the difference lies between
\(-\epsilon^{-1}\Delta(\mathbf w,\mathbf w')\) and
\((1+\epsilon^{-1})\Delta(\mathbf w,\mathbf w')\).
Interchanging the vectors covers the remaining mixed regime.

Finally, \(x\mapsto[x]_+\) is one-Lipschitz. Applying this to the hinge in \cref{obj:def:dual-process} gives
\[
|f_\lambda(\mathbf w)-f_\lambda(\mathbf w')|
\leq
2|g_0(\mathbf w)-g_0(\mathbf w')|
+|g_1(\mathbf w)-g_1(\mathbf w')|
+|\lambda|\,|s(\mathbf w)-s(\mathbf w')|
\leq\Lambda_\epsilon\Delta(\mathbf w,\mathbf w').
\]
Since \(\Delta(\mathbf w,\mathbf w')\leq4\|\mathbf w-\mathbf w'\|_\infty\), \(f_\lambda\) is continuous on \(\mathcal R\). % lean: CausalSmith.Stat.DiscreteBudgetvalueCurve.budget_thresholdFunReal_pointwise_lipschitz

\item Write \(\beta_\alpha^{c^0,R}(\lambda)\) for the centered coefficients in \cref{obj:def:jf-estimator} with its pilot center and radii replaced by \(c^0,R\). Explicitly, they are the unique coefficients satisfying
\[
P_\lambda(c^0+R\odot\mathbf z)-f_\lambda(c^0)
=
\sum_{\alpha\in\{0,\ldots,2(K-1)\}^{\mathcal Z}}
\beta_\alpha^{c^0,R}(\lambda)
\prod_{\zeta\in\mathcal Z}z_\zeta^{\alpha_\zeta},
\qquad \mathbf z\in\mathbb R^{\mathcal Z}.
\]
The normalized polynomial is therefore
\[
\widetilde P_\lambda(\mathbf z)
=
f_\lambda(c^0)+
\sum_{\alpha\in\{0,\ldots,2(K-1)\}^{\mathcal Z}}
\beta_\alpha^{c^0,R}(\lambda)
\prod_{\zeta\in\mathcal Z}z_\zeta^{\alpha_\zeta},
\qquad \mathbf z\in\mathbb R^{\mathcal Z}.
\]
The continuity just proved applies to
\(\mathbf z\mapsto f_\lambda(c^0+R\odot\mathbf z)\) on \([-1,1]^4\).
The normalized construction, including the added center value, therefore gives, for every \(\boldsymbol\vartheta\in\mathbb R^4\),
\[
\widetilde P_\lambda(\cos\boldsymbol\vartheta)
=
\int_{[-\pi,\pi]^4}
f_\lambda\!\left(c^0+R\odot
\cos(\boldsymbol\vartheta+\boldsymbol\omega)\right)
\prod_{\zeta\in\mathcal Z}J_K(\omega_\zeta)
\,d\boldsymbol\omega.
\]
The order-four kernel has trigonometric degree \(2(K-1)\) in each coordinate, so \(\widetilde P_\lambda\) has coordinatewise polynomial degree at most \(2(K-1)\). Positive radii permit the affine substitution
\[
P_\lambda(\mathbf w)
=
\widetilde P_\lambda\!\left(
\left(\frac{w_\zeta-c^0_\zeta}{R_\zeta}\right)_{\zeta\in\mathcal Z}
\right),
\qquad \mathbf w\in\mathbb R^4,
\]
which preserves those coordinatewise degree bounds and identifies this polynomial with the one in the statement. % lean: Causalean.Mathlib.Analysis.JacksonApproximation.mvPolynomial_affine_substitution_four

\item Put
\[
\vartheta_\zeta
=
\arccos\!\left(\frac{y_\zeta-c^0_\zeta}{R_\zeta}\right),
\qquad \zeta\in\mathcal Z.
\]
The evaluation-point assumption makes every angle well defined in \([0,\pi]\), and \(y=c^0+R\odot\cos\boldsymbol\vartheta\). Positivity and unit integral of each \(J_K\), followed by the Lipschitz estimate from the first step, yield
\[
|P_\lambda(y)-f_\lambda(y)|
\leq
\Lambda_\epsilon\sum_{\zeta\in\mathcal Z}R_\zeta
\int_{[-\pi,\pi]^4}
\left|\cos(\vartheta_\zeta+\omega_\zeta)
-\cos\vartheta_\zeta\right|
\prod_{\xi\in\mathcal Z}J_K(\omega_\xi)
\,d\boldsymbol\omega.
\]
The cosine addition formula and
\(|\sin t|\leq|t|\), \(1-\cos t\leq t^2/2\) give
\[
|\cos(\vartheta_\zeta+\omega_\zeta)-\cos\vartheta_\zeta|
\leq|\sin\vartheta_\zeta|\,|\omega_\zeta|
+\frac{\omega_\zeta^2}{2}.
\]

For completeness, the kernel moments used here obey
\[
\int_{-\pi}^{\pi}|t|J_K(t)\,dt\leq\frac{32}{K},
\qquad
\int_{-\pi}^{\pi}t^2J_K(t)\,dt\leq\frac{64}{K^2}.
\]
Indeed, writing \(D_K(t)=\sum_{r=0}^{K-1}e^{irt}\), the unnormalized fourth power in \cref{obj:def:jf-estimator} is \(|D_K(t)|^4\). Orthogonality of the Fourier modes gives
\[
\int_{-\pi}^{\pi}|D_K(t)|^4\,dt
=2\pi\sum_{h=-(K-1)}^{K-1}(K-|h|)^2
=\frac{2\pi}{3}K(2K^2+1)
\geq\frac{4\pi}{3}K^3.
\]
For \(t\in[-\pi,\pi]\), the inequality
\(t^2\leq\pi^2\sin^2(t/2)\), together with
\(\int_{-\pi}^{\pi}\{\sin(Kt/2)/\sin(t/2)\}^2dt=2\pi K\),
bounds its weighted second integral by \(2\pi^3K\).
Normalization and \(\pi^2\leq16\) thus give the sharper second-moment bound \(24/K^2\). Integrating
\[
|t|\leq\frac K4t^2+\frac1K
\]
gives a first-moment bound \(7/K\); the displayed bounds follow for every positive integer \(K\).

Integrating the cosine estimate coordinatewise now gives
\[
|P_\lambda(y)-f_\lambda(y)|
\leq
32\Lambda_\epsilon
\sum_{\zeta\in\mathcal Z}
\left\{
\frac{R_\zeta|\sin\vartheta_\zeta|}{K}
+\frac{R_\zeta}{K^2}
\right\}.
\]
Finally,
\[
R_\zeta|\sin\vartheta_\zeta|
=
\sqrt{(y_\zeta-c^0_\zeta+R_\zeta)
(c^0_\zeta+R_\zeta-y_\zeta)}.
\]
Substitution proves the asserted bound, including points on the faces of the rectangle. % lean: Causalean.Mathlib.Analysis.Approximation.Chebyshev.AffineFourBoundary.affineJackson_boundary_adaptive_four
\end{enumerate}
\end{proof}
